% part: intuitionistic-logic
% chapter: propositions-as-types
% section: introduction

\documentclass[../../../include/open-logic-section]{subfiles}

\begin{document}

\olfileid{int}{pty}{int}

\olsection{Pambuka}

Kalkulus lambda iku modhèl komputasi kang prasaja.
Kalkulus iki ngolah tèrma kang dibangun saka variabel.
Yèn $N$ lan $M$ iku tèrma, $(NM)$ uga tèrma
(``$N$ ditrapaké marang~$M$''). Yèn $N$ iku tèrma,
$\lambd[x][N]$ uga tèrma. Yèn $N$ ngemot variabel~$x$,
kedadéan~$x$ kang cocog ing $\lambd[x][N]$ kaiket
déning operator $\lambd[x]$, mula ora bébas manèh.
Tèrma awujud $(\lambd[x][N]M)$ diarani ngalami
kontraksi~$\beta$ dadi $\Subst{N}{M}{x}$
(asil ngganti saben kedadéan bébas $x$ ing~$N$
nganggo $M$, kanthi ngganti jeneng variabel kaiket
luwih dhisik yèn perlu supaya variabel bébas ora
kaiket). Tèrma $N$ ngalami konversi~$\beta$
dadi~$N'$ yèn $N'$ diasilaké saka $N$ kanthi
kontraksi~$\beta$ ing sub-tèrma; kita nulis
$N \redone N'$. Komputasi ing kalkulus~$\lambd$
iku rerangkéyan $N \redone N_1 \redone N_2 \redone \dots$.
Yèn rerangkéyan iki pungkasané ing tèrma $N_k$
kang ora nduwèni sub-tèrma kang bisa ngalami
kontraksi~$\beta$, tèrma mau diarani \emph{normal}
utawa wujud normal saka~$N$. Komputasi ing
kalkulus~$\lambd$ bisa dianggep rerangkéyan kaya
mangkono. Komputasi mandheg yèn ngasilaké wujud
normal, lan wujud normal iku asil komputasiné.
Modhèl prasaja iki jebul lengkap miturut Turing:
saben fungsi rékursif parsial bisa dikomputasi
nganggo tèrma ing kalkulus~$\lambd$.

Haskell Curry lan William Howard kanthi mandhiri
nemokaké pepadhan kang raket antarané deduksi alamiah
lan kalkulus~$\lambd$. Miturut intuisi, tèrma
$\lambd[x][N]$ iku fungsi: $x$ argumèné, lan $N$
nuduhaké cara ngétung bijiné yèn diwènèhi~$x$.
Tèrma $(\lambd[x][N]M)$ banjur ngetrapaké fungsi
$\lambd[x][N]$ marang argumèn~$M$; kontraksi~$\beta$
nuduhaké cara évaluasi: ganti $x$ ing $N$ nganggo~$M$
(banjur terus évaluasi tèrma asilé). Saiki anggep
iki fungsi kanthi domain lan jangkauan tetep.
Yèn domain $\lambd[x][N]$ iku $A$ lan jangkauané~$B$,
variabel~$x$ mung nampa argumèn saka~$A$, lan $N$
ngasilaké biji saka~$B$. Mula $(\lambd[x][N]M)$ mung
nduwèni teges yèn $\lambd[x][N]$ fungsi $A \to B$
lan $M \in A$. Kanthi nambah katerangan iki ing
kalkulus~$\lambd$, kita olèh kalkulus \emph{mawa tipe}
$\lambd$. Ing kalkulus~$\lambd$ mawa tipe,
ora saben èksprèsi $(\lambd[x][N]M)$ diidinaké,
nanging mung yèn ``tipe'' saka $\lambd[x][N]$
lan $M$ cocog: $\lambd[x][N]$ nduwèni tipe
$A \lif B$, lan $M$ nduwèni tipe~$A$.
Tipe $\lambd[x][N]$ iku $A \to B$ mung yèn
tipe $x$ iku $A$ lan tipe $N$ iku~$B$.
Umumé, ora saben èksprèsi $(M_1M_2)$ diidinaké.
Mung tèrma $(M_1M_2)$ kang $M_1$ nduwèni tipe
$A \to B$ lan $M_2$ nduwèni tipe~$A$ kang
diidinaké; ing kéné $(M_1M_2)$ nduwèni tipe~$B$.

Aturan tipe iki cetha cocog karo !!{derivation}s
ing deduksi alamiah: tipe diwakili !!{formula}s,
lan !!{derivation}s cocog karo tèrma~$\lambd$.
Tuladhané, !!a{derivation} tumrap $!A \lif !B$
cocog karo tèrma $\lambd[\typeof{x}{!A}][N]$,
kang nduwèni tipe fungsi $!A \lif !B$.
Aturan inferènsi deduksi alamiah cocog karo aturan
tipe ing kalkulus lambda mawa tipe, kayata:
\begin{prooftree}
  \AxiomC{$\Discharge{!A}{x}$}
  \DeduceC{$!B$}
  \DischargeRule{\Intro\lif}{x}
  \UnaryInfC{$!A \lif !B$}
  \DisplayProof\bottomAlignProof
  \quad\text{cocog karo}\quad
  \Axiom$x: !A \fCenter N: !B$
  \UnaryInf$ \fCenter \lambd[\typeof{x}{!A}][N] : !A \lif !B$
\end{prooftree}
Aturan ing sisih tengen ateges yèn $x$ nduwèni
tipe~$!A$ lan $N$ nduwèni tipe~$!B$,
mula $\lambd[\typeof{x}{!A}][N]$ nduwèni
tipe~$!A \lif !B$.

Aturan $\Elim\lif$ cocog karo aturan tipe kanggo
tèrma aplikasi, yaiku:
\[
  \AxiomC{$!A \lif !B$}
  \AxiomC{$!A$}
  \RightLabel{\Elim\lif}
  \BinaryInfC{$!B$}
  \DisplayProof
  \quad\text{cocog karo}\quad
  \Axiom$\fCenter M_1: !A \lif !B$
  \Axiom$\fCenter M_2: !A$
  \BinaryInf$ \fCenter (M_1M_2) : !B$
  \DisplayProof
\]
Yèn aturan \Intro\lif{} langsung disusul aturan
\Elim\lif{}, !!{derivation} bisa diprasajakaké:
\begin{gather*}
  \AxiomC{$\Discharge{!A}{x}$}
  \DeduceC{$!B$}
  \DischargeRule{\Intro\lif}{x}
  \UnaryInfC{$!A \lif !B$}
  \AxiomC{}
  \DeduceC{$!A$}
  \RightLabel{\Elim\lif}
  \BinaryInfC{$!B$}
  \DisplayProof
  \quad\redone\quad
  \AxiomC{}
  \DeduceC{$!A$}
  \DeduceC{$!B$}
  \DisplayProof
\end{gather*}
kang cocog karo kontraksi~$\beta$ ing tèrma lambda:
\[
(\lambd[\typeof{x}{!A}][N]M) \quad\redone\quad
\Subst{N}{M}{x}.
\]

Pepadhan kang padha uga lumaku antarané aturan
$\land$ lan tipe ``prodhuk'', sarta antarané
aturan $\lor$ lan tipe ``gunggung''.

Sesambungan antarané tèrma ing kalkulus lambda
mawa tipe prasaja lan !!{derivation}s deduksi
alamiah iki diarani sesambungan ``Curry--Howard'',
``bukti minangka program'', utawa ``proposisi
minangka tipe''. Saliyané !!{formula}s (proposisi)
kang cocog karo tipe lan bukti kang cocog karo
tèrma, sesambungan-sesambungan iki bisa diringkes
mangkéné:
\begin{center}
  \begin{tabular}{c c c}
    logika & program \\
    \hline
    proposisi/!!{formula} & tipe \\
    bukti & tèrma \\
    asumsi & variabel \\
    asumsi kang dibusak & variabel kaiket \\
    asumsi kang isih mbukak & variabel bébas \\
    $!A \lif !B$ & tipe fungsi \\
    $!A \land !B$ & tipe prodhuk \\
    $!A \lor !B$ & tipe gunggung \\
    $\lfalse$ & tipe kosong \\
    \hline
  \end{tabular}
\end{center}

Sesambungan Curry--Howard dadi salah siji dhasar
asistèn bukti otomatis lan pamriksa tipe program,
amarga mriksa bukti kang nyeksèni proposisi
(kaya ing ndhuwur) padha karo mriksa apa program
(tèrma) nduwèni tipe kang diumumaké.
\end{document}
